\hypertarget{appendix-40df0855b4a7}{}\label{referencestart:appendix-40df0855b4a7}
\section{Automatic differentiation and backpropagation}\label{appendix-40df0855b4a7-automatic-differentiation-and-backpropagation}

\textbf{Curriculum depth:} Developed · \textbf{First seen:} University courses · \textbf{Latest development:} Phase 9

\subsection{Definition and mechanism}\label{appendix-40df0855b4a7-definition-and-mechanism}

Automatic differentiation evaluates derivatives by applying the chain rule to the operations in a program. Each operation has a local derivative rule; composing the rules gives the derivative of the complete computation without symbolically expanding it or approximating it by finite differences. Forward mode propagates a tangent perturbation and computes a Jacobian-vector product (JVP), useful when there are few inputs. Reverse mode propagates sensitivities from outputs back to inputs and computes a vector-Jacobian product (VJP), making gradients of one scalar loss with respect to millions of parameters practical. Backpropagation is reverse-mode differentiation through a neural network's computation graph. It does not remove numerical problems: nondifferentiable operations require conventions, saturated activations can yield tiny gradients, and an incorrect graph or reduction yields an incorrect gradient. Gradients are local sensitivity information, not proof of a global optimum. Finite differences remain a useful small-scale check, but cancellation and step-size tradeoffs make them a poor training method.

\subsubsection{Phase 9--10 development}\label{appendix-40df0855b4a7-phase-910-development}

Forward mode propagates an input tangent through local derivatives and computes a Jacobian-vector product. Reverse mode propagates an output sensitivity backward and computes a vector-Jacobian product. These are chain-rule evaluations on the operation graph, distinct from finite differences, which perturb inputs and can suffer cancellation. The reverse pass requires a valid reverse-topological order and storage that remains alive until propagation finishes. Differentiating a converged solution uses an implicit equation and its local Jacobian; it is not automatically the same as differentiating every solver iteration.

\subsection{Prerequisites}\label{appendix-40df0855b4a7-prerequisites}

\begin{itemize}
\tightlist
\item
  \hyperref[appendix-29f4bea4c2bb]{Constrained optimization and inverse kinematics}
\item
  \href{../07-Vision-XR-and-Playable-Worlds/main.pdf\#nameddest=appendix-ab31ef447dbf}{Matrix multiplication and composition order (Vision, XR, and Playable Worlds, Reference notes: Matrix multiplication and composition order)}
\item
  \href{../01-Bits-Programs-and-Reliable-Systems/main.pdf\#nameddest=appendix-c2ac4a84c283}{Floating-point reassociation (Bits, Programs, and Reliable Systems, Reference notes: Floating-point reassociation)}
\item
  \href{../07-Vision-XR-and-Playable-Worlds/main.pdf\#nameddest=appendix-7bc181935c42}{Derivative as a local linear map (Vision, XR, and Playable Worlds, Reference notes: Derivative as a local linear map)}
\item
  \href{../07-Vision-XR-and-Playable-Worlds/main.pdf\#nameddest=appendix-d8a7ddb675eb}{Jacobians, pushforwards, and pullbacks (Vision, XR, and Playable Worlds, Reference notes: Jacobians, pushforwards, and pullbacks)}
\end{itemize}

\subsection{Used by}\label{appendix-40df0855b4a7-used-by}

\begin{itemize}
\tightlist
\item
  \href{../07-Vision-XR-and-Playable-Worlds/main.pdf\#nameddest=04-visual-computing-ii-ch031-learning-a-visual-classifier-from-examples}{Learning a visual classifier from examples (Vision, XR, and Playable Worlds, Vision, XR, and Playable Worlds, Chapter 2)}
\item
  \href{../05-Spatial-Worlds-Engine-State-and-Simulation/main.pdf\#nameddest=04-visual-computing-ii-ch029-optimization-constraints-and-inverse-kinematics}{Optimization, constraints, and inverse kinematics (Spatial Worlds, Engine State, and Simulation, Spatial Worlds, Engine State, and Simulation, Chapter 10)}
\item
  \hyperref[appendix-47c51ea026fb]{Learned image features and CNNs}
\item
  \hyperref[appendix-a9d5dfdd164a]{Attention and visual transformers}
\item
  \hyperref[appendix-2c568f05ee34]{Latent-variable models and generative flows}
\item
  \hyperref[appendix-c5b332e400e7]{Radiance fields and Gaussian splatting}
\item
  \hyperref[appendix-b0c42da325f1]{Perception feedback and robot learning}
\item
  \href{../07-Vision-XR-and-Playable-Worlds/main.pdf\#nameddest=appendix-8836553d56fc}{Implicit differentiation and parameter sensitivities (Vision, XR, and Playable Worlds, Reference notes: Implicit differentiation and parameter sensitivities)}
\end{itemize}

\subsection{Full teaching explanation}\label{appendix-40df0855b4a7-full-teaching-explanation}

\begin{itemize}
\tightlist
\item
  \href{../07-Vision-XR-and-Playable-Worlds/main.pdf\#nameddest=04-visual-computing-ii-ch031-learning-a-visual-classifier-from-examples}{Learning a visual classifier from examples (Vision, XR, and Playable Worlds, Vision, XR, and Playable Worlds, Chapter 2)}
\item
  \href{../05-Spatial-Worlds-Engine-State-and-Simulation/main.pdf\#nameddest=04-visual-computing-ii-ch029-optimization-constraints-and-inverse-kinematics}{Optimization, constraints, and inverse kinematics (Spatial Worlds, Engine State, and Simulation, Spatial Worlds, Engine State, and Simulation, Chapter 10)}
\item
  \href{../02-Mathematics-for-Visual-Computing/main.pdf\#nameddest=02-foundations-ii-ch022-numerical-differentiation-and-automatic-differentiation}{Phase 9 teaching section (Mathematics for Visual Computing, Mathematics for Visual Computing, Chapter 6)}
\item
  \href{../01-Bits-Programs-and-Reliable-Systems/main.pdf\#nameddest=appendix-adf48f5dc5c1}{Visual Recall/09 - Calculus and Optimization - Visual Recall.md (Bits, Programs, and Reliable Systems, Reference notes: 09 - Calculus and Optimization - Visual Recall)} · Phase 9 visual recall
\end{itemize}

\subsection{Recall and next depth}\label{appendix-40df0855b4a7-recall-and-next-depth}

\begin{itemize}
\tightlist
\item
  Compare a JVP and VJP by naming the direction propagated and the Jacobian multiplication performed.
\item
  Why does a correct gradient still fail to guarantee a global minimum?
\end{itemize}

\subsection{Source and navigation}\label{appendix-40df0855b4a7-source-and-navigation}

Teaching narrative: \href{../07-Vision-XR-and-Playable-Worlds/main.pdf\#nameddest=04-visual-computing-ii-ch031-learning-a-visual-classifier-from-examples}{open the main explanation (Vision, XR, and Playable Worlds, Vision, XR, and Playable Worlds, Chapter 2)}. Browse the \hyperref[appendix-a6d556a4ef50-university-graphics-vision-and-photography---independent-series]{University graphics and vision index}. Additional chapter connections: \href{../05-Spatial-Worlds-Engine-State-and-Simulation/main.pdf\#nameddest=04-visual-computing-ii-ch029-optimization-constraints-and-inverse-kinematics}{Optimization, constraints, and inverse kinematics (Spatial Worlds, Engine State, and Simulation, Spatial Worlds, Engine State, and Simulation, Chapter 10)}; \href{../02-Mathematics-for-Visual-Computing/main.pdf\#nameddest=01-foundations-i-ch020-qr-turns-least-squares-into-a-triangular-solve-plus-an-irreducible-residual}{QR turns least squares into a triangular solve plus an irreducible residual (Mathematics for Visual Computing, Mathematics for Visual Computing, Chapter 4)}; \href{../01-Bits-Programs-and-Reliable-Systems/main.pdf\#nameddest=01-foundations-i-ch002-a-compiled-interface-rather-than-repeated-text}{A compiled interface rather than repeated text (Bits, Programs, and Reliable Systems, Bits, Programs, and Reliable Systems, Chapter 2)}.

Caption/source records: \href{../07-Vision-XR-and-Playable-Worlds/main.pdf\#nameddest=04-visual-computing-ii-ch036-sc-02---stanford-cs231n-spring-2025-lecture-2-image-classification-with-linear-classifiers}{SC-02 (Vision, XR, and Playable Worlds, Vision, XR, and Playable Worlds, Chapter 5)}, \href{../07-Vision-XR-and-Playable-Worlds/main.pdf\#nameddest=04-visual-computing-ii-ch036-sc-03---stanford-cs231n-spring-2025-lecture-3-regularization-and-optimization}{SC-03 (Vision, XR, and Playable Worlds, Vision, XR, and Playable Worlds, Chapter 5)}, \href{../07-Vision-XR-and-Playable-Worlds/main.pdf\#nameddest=04-visual-computing-ii-ch036-sc-04---stanford-cs231n-spring-2025-lecture-4-neural-networks-and-backpropagation}{SC-04 (Vision, XR, and Playable Worlds, Vision, XR, and Playable Worlds, Chapter 5)}.
